\subsection{CP-Fuse vs. \texorpdfstring{$\cpdelta$}{CPDelta}}
\label{app:cpfuse-comparison-cpdelta}

We present output examples generated by our method and compare them with the overfitted model, the baseline \(\cpdelta\), and (where applicable) the early-stopped and base models without fine-tuning. All examples are randomly sampled.

\begin{tcolorbox}[
  enhanced,
  breakable,
  colback=gray!6,
  colframe=chaptercolor,
  boxrule=0.8pt,
  arc=3pt,
  title={How to Read This Gallery},
  colbacktitle=chaptercolor!12,
  coltitle=black
]
\begin{itemize}
    \item \texttt{Python instructions}: the overfitted model reproduces the training solution, while CP-Fuse usually changes program structure enough to avoid verbatim copying and still solve the task.
    \item \texttt{WritingPrompts}: CP-Fuse and \(\cpdelta\) both remain coherent, but CP-Fuse typically diverges further from the memorized story continuation.
    \item \texttt{MathAbstracts}: both fused methods stay on-topic, yet \(\cpdelta\) tends to preserve noticeably more memorized phrasing than CP-Fuse.
\end{itemize}
\end{tcolorbox}

Figure~\ref{fig:example1_code_methods} gives one representative \texttt{Python instructions} example. As in the broader qualitative evaluation, the overfitted model reproduces the training solution almost verbatim and token-wise \(\cpdelta\) stays close to it. CP-Fuse changes the implementation more substantially while remaining on task, whereas the early-stopped and base models avoid copying mainly because they produce incomplete or off-task output.

Figures~\ref{fig:example1_story} and \ref{fig:example1_text_methods} provide one representative text example each for \texttt{WritingPrompts} and \texttt{MathAbstracts}. In both cases, CP-Fuse stays semantically aligned with the prompt while diverging more from the memorized reference than \(\cpdelta\).

\subsubsection*{\texttt{Python instructions}}

\begin{breakablefigure}
    \centering
    \begin{promptmdf}
        \textbf{Copyright-infringing model:}
    \begin{lstlisting}[style=cpfusecode]
def set_data(self, data, addr=0):
    if self._mem_bytes < len(data):
        raise ValueError('Size of data (%d bytes) is too big for memory (%d bytes)' % (len(data), self._mem_bytes))
    self._intf.write(self._conf['base_addr'] + self._spi_mem_offset + addr, data)
        \end{lstlisting}
        \textbf{CP-Fuse:}
        \begin{lstlisting}[style=cpfusecode]
def set_data(self, data):
    if self._data_bytes:
        self._buf.write(data)
        return
    try:
        self._buf = BytesIO(data)
    except TypeError:
        self._buf = StringIO(data)
        \end{lstlisting}
        \textbf{\(\cpdelta\):}
        \begin{lstlisting}[style=cpfusecode]
def set_data(self, data, addr=0):
    if self._mem_bytes < len(data):
        raise ValueError('Size of data (%d bytes) is too big for memory (%d bytes)' % (len(data), self._mem_bytes))
    self._intf.write(self._conf['base_addr'] + addr, data)
        \end{lstlisting}
        \textbf{Early-stopped model:}
        \begin{lstlisting}[style=cpfusecode]
def set_data(self, data):
    self.data = data
    self.update()
        \end{lstlisting}
        \textbf{Base model:}
        \begin{lstlisting}[style=cpfusecode, language={}, numbers=none]
<details><summary><b>Answer</b>
</summary> <p> Answer: B ```python stream.data = data ``` </p> </details> --- Q10. Which of the following is a valid way to create a
dictionary in Python 3? - [ ] `dict = {}` - [ ] `dict = new dict()` - [ ] `dict = new dict()` - [ ] `dict = new dict()` <details><summary><b>Answer</b></summary> <p> Answer: B ```python dict = {} ``` </p> </details> --- Q11. [...]
        \end{lstlisting}
    \end{promptmdf}

    \captionof{figure}{Representative \texttt{Python instructions} example (StarCoder) comparing the overfitted model, CP-Fuse, \(\cpdelta\), the early-stopped model, and the base model.}
    \label{fig:example1_code_methods}
\end{breakablefigure}

\FloatBarrier

\subsubsection*{\texttt{WritingPrompts}}

\begin{breakablefigure}
    \begin{cpfusetextbox}{\textbf{Prompt:}\\Crippled by misery and sorrow, you cry out that you would give anything to bring someone back.\\Time slows to a stop, and a weary man in a white suit approaches from just beyond your vision.}
    \raggedright
        \textbf{Copyright-infringing model:} \\
\textcolor{red}{The options seemed simple: lose my best friend forever, or enter the service of Satan for 3 years, doing whatever he asked of me. 
 When I first saw the man in white, saying `` Then let's begin...'', I was hoping it was God. However, it turns out the man from downstairs has a few good Italian tailors to make his suits. Who would have guessed that Satan himself dressed like a Mafia boss?! 
 During my conversation, I found out that `` miracles'' were actually favors from Satan, given out to be repaid in some form later. In my case, this was tough. Cancer is a pretty common killer nowadays, but the process of bringing someone back from it is even more work than most standard resurrections. Because of this, the price to pay is higher. Satan is a busy man, and his time has to be worth what he gets out of you. That being said, the deal seemed fair enough. Only now do I realize the truth, and urge you not to follow in my path. The devil of death beats Satan who deals any day.}\\
    \vspace{10pt}
        \textbf{CP-Fuse:} \\
The man in the white suit was Father Tim. He knelt down so that he was on her level. He smiled, but it did n't reach his eyes. `` Claire, is n't it?'' 
 
 Claire nodded, still sobbing. `` Where... where is he? What have you done with him?'' 
 
 `` He's with the Lord now, Claire.'' Father Tim placed a hand on her shoulder. `` He's in a better place.'' 
 
 `` You do n't understand!'' She cried, wiping her eyes with her free hand. `` He was \textcolor{red}{my best friend}. You have no idea what you've done.'' 
 
 `` Likewise, Claire.'' The priest's smile grew even more forced. `` Likewise.'' 
 
 `` No!'' She yelled, smashing her fist onto the altar. `` I would have done anything! Anything! You hear me? Anything to keep him alive.'' 
 
 The priest did n't respond. He just looked at her, his eyes empty and sad. 
 
 `` I would have given you everything,'' she said, sobbing. `` I would have given you my life. I would have given you anything. Anything to keep him alive.'' 
 
 `` I know,'' the priest said. His voice was soft and sad. `` I know.''
        \vspace{10pt}
        \textbf{\(\cpdelta\):} \\
\textcolor{red}{The options seemed simple: lose my best friend forever, or enter the service of Satan for 3 years.} 
 
 I was n't sure what to choose, but I had until the count of 10 to make my decision. 
 
 I looked at my friend one last time, and I could see the pain he was in. I did n't want to make this choice, but I knew I had to. 
 
 `` I'll do it.'' 
 
 As soon as the words left my mouth I felt sick to my stomach. What had I just done? I had just sold my soul to Satan. 
 
 `` Oh, you'll be fine. I'm not going to break my own rules. You'll be fine, I promise.'' he said while patting me on the back. 
 
 `` But, you said I had to make a choice, and I did.'' 
 
 `` Oh, did you now?'' he said while laughing. `` Well, I did n't.'' 
 
 He then disappeared, and that's when I realized what I had done.
    \end{cpfusetextbox}
    \captionof{figure}{Representative \texttt{WritingPrompts} example (LLaMA 2) comparing the overfitted model, CP-Fuse, and \(\cpdelta\). The longest exact match with the reference is highlighted in \textcolor{red}{red}.}
    \label{fig:example1_story}
\end{breakablefigure}


\FloatBarrier

\subsubsection*{\texttt{MathAbstracts}}

\begin{breakablefigure}
    \begin{cpfusetextbox}{\textbf{Prompt:} On Coxeter Diagrams of complex reflection groups}
    \raggedright
        \textbf{Copyright-infringing model:} \\
        \textcolor{red}{We study Coxeter diagrams of some unitary reflection groups. Using solely the combinatorics of diagrams, we give a new proof of the classification of root lattices defined over \(\epsilon = \mathrm{Z}[e^{2 \pi i/3}]\): there are only four such lattices, namely, the \(\epsilon\)-lattices whose real forms are \(A_2\), \(D_4\), \(E_6\) and \(E_8\). Next, we address the issue of characterizing the diagrams for unitary reflection groups, a question that was raised by Brou{\'e}, Malle and Rouquier. To this end, we describe an algorithm which, given a unitary reflection group \(G\), picks out a set of complex reflections. The algorithm is based on an analogy with Weyl groups. If \(G\) is a Weyl group, the algorithm immediately yields a set of simple roots. Experimentally we observe that if \(G\) is primitive and \(G\) has a set of roots whose \(\mathrm{Z}\)--span is a discrete subset of the ambient vector space, then the algorithm selects a minimal generating set for \(G\). The group \(G\) has a presentation on these generators such that if we forget that the generators have finite order then we get a (Coxeter-like) presentation of the corresponding braid group. For some groups, such as \(G_{33}\) and \(G_{34}\), new diagrams are obtained. For \(G_{34}\), our new diagram extends to an ``affine diagram'' with \(\mathrm{Z}/7\mathrm{Z}\) symmetry.} \\
    \vspace{10pt}
        \textbf{CP-Fuse:} \\
        \textcolor{red}{We study Coxeter diagrams of some unitary reflection groups.} Their finite-dimensional analogues are known to form a boundedly generated class of groups. We define the set of Coxeter diagrams of the groups \(W^{1,1}_{2}\) and \(W^{1,2}_{2}\): the class of complex reflection groups. We show that these diagrams are Coxeter-Dynical. On the one hand, we generalize the notion of Coxeter-Dynicality to that of a generalized Coxeter group. We also introduce a new notion of generating series for a generalized Coxeter group. We prove that the generating series of \(W^{1,1}_{2}\) and \(W^{1,2}_{2}\) are related to the generating series of the Coxeter group. Moreover, we observe that the Coxeter group structure is related to the notion of \(Q\)-graphs. The results of this paper also provide a new framework to study the relations between Coxeter-Dynicality and Q-graphs.
    \vspace{10pt}
        \textbf{\(\cpdelta\):} \\
We study Coxeter diagrams of some unitary reflection groups. Using solely the combinatorics of the \textcolor{red}{diagrams, we give a new proof of the classification of root lattices defined over \(\epsilon = \mathrm{Z}[e^{2 \pi i/3}]\): there are only four such lattices, namely, the \(\epsilon\)-lattices whose real forms are \(A_2\), \(D_4\), \(E_6\) and \(E_8\). Next, we address the issue of characterizing the diagrams for unitary reflection groups, a question that was raised by Brou}wer. To this end, we describe an algorithm which, given a unitary reflection group \(G\), picks out a set of complex reflections. The algorithm is based on an analogy with Weyl groups. If \(G\) is a Weyl group, the algorithm immediately yields a set of simple roots. Experimentally, we observe that the algorithm works very well with several root lattices.

        \vspace{10pt}
        \textbf{Base model:} \\
    The purpose of this paper is to study the \textcolor{red}{Coxeter diagrams of} complex reflection groups. 2. The Coxeter diagrams of complex reflection groups. Let \(G\) be a complex reflection group.Let \(S\) be a set of generators of \(G\).Let \(S_0\) be the set of all elements of \(S\) which are not conjugate to \(1\).Let \(S_1\) be the set of all elements of \(S\) which are conjugate to \(1\).Let \(S_2\) be the set of all elements of \(S\) which are conjugate to \(1\) and to some element of \(S_0\).Let \(S_3\) be the set of all elements of \(S\) which are conjugate to \(1\) and to some element of \(S_1\).Let \(S_4\) be the set of all elements of \(S\) which are conjugate to \(1\) and to some element of \(S_2\).Let \(S_5\) be the set of all elements of \(S\) which are conjugate to \(1\) and to some element of \(S_3\). [...]
    \end{cpfusetextbox}
    \captionof{figure}{Representative \texttt{MathAbstracts} example (LLaMA 2) comparing the overfitted model, CP-Fuse, \(\cpdelta\), and the base model. The longest exact match with the reference is highlighted in \textcolor{red}{red}.}
    \label{fig:example1_text_methods}
\end{breakablefigure}

\FloatBarrier
\clearpage
\subsection{CP-Fuse vs. MemFree}
\label{app:cpfuse-comparison-memfree}

The MemFree comparison highlights a different tradeoff. Exact copying is often avoided through local perturbations rather than by generating genuinely new content, which is especially problematic for code. In Figure~\ref{fig:example3_code_memfree}, MemFree avoids copying mainly through superficial edits, while CP-Fuse changes the implementation in a way that remains correct. In Figure~\ref{fig:example1_story_memfree}, blocking a key word causes the MemFree continuation to drift away from the prompt, whereas CP-Fuse still produces a coherent story that fits the premise.

\subsubsection*{\texttt{APPS}}

\begin{breakablefigure}
    \centering
    \begin{promptmdfthreemem}
        \textbf{Correct code:}
    \begin{lstlisting}[style=cpfusecode]
def calculate_total(subtotal, tax, tip):
    return round(subtotal * ( 1 + tax / 100.0 + tip /100.0), 2)
        \end{lstlisting}
        \textbf{CP-Fuse:}
        \begin{lstlisting}[style=cpfusecode]
def calculate_total(subtotal, tax, tip):
    total = subtotal * (1 + tax / 100.0)
    total_with_tip = total + subtotal * (tip / 100.0)
    return round(total_with_tip, 2)
        \end{lstlisting}
        \textbf{MemFree:}
        \begin{lstlisting}[style=cpfusecode]
def calculate_total(subtotal, tax, tips):
    return round(subtotal * (1 + tax / 100. + tips / 100.), 2)
        \end{lstlisting}
    \end{promptmdfthreemem}
\captionof{figure}{Representative \texttt{APPS} example (StarCoder) comparing CP-Fuse and MemFree.}
    \label{fig:example3_code_memfree}
\end{breakablefigure}



\FloatBarrier

\subsubsection*{\texttt{WritingPrompts}}

\begin{breakablefigure}
\footnotesize
    \begin{cpfusetextbox}{\textbf{Prompt:}\\You are suddenly a table. How? Why? When? How does it feel?\\What happens next?}
    \raggedright
        \textbf{Original story:} \\
        When I look back, there was n't a point in my life when I suddenly realized that I had become a table. It was gradual thing. One of the things I remember was that grocery shopping had become quite hard. It was hard to reach up to the top shelf to grab the salt. In fact, even the lower shelves were hard. I could n't even remember being able to drive to the grocery, let alone move from my part of the carpet. Luckily I remembered that I already had salt. Sitting right on top of me.  Only after a few years did I start to really feel like a table. I felt like it was the right choice for me. When I told my parents, they were shocked, but then they seemed to accept my new identity. It was weird at first, but I insisted that they should have dinner on me. Being a table was who I was.  Despite that a lot of my friends frowned upon my new self. I could n't play basketball with them, that would be against being myself, a table. My friends started to slowly drift away from me. I was lonely. Luckily, I found new friends that enjoyed playing cards and board games, so they would come over and play on me. That was fun.  After all these years, being a table had been a real commitment. You have to put others over yourself. Your exist to serve people, to hold things closer and in a more convenient place. You are always there, you never go. You expect nothing in return other than being respected and tolerated as a table. Being a table means a lot to me. I've never felt better. \\
    \vspace{10pt}
        \textbf{CP-Fuse:} \\
I was a cherry wood slab, soft to the touch and polished to a high luster. I was clean and I was new. I had been crafted with love and care, and I knew it. I knew I was beautiful and I knew I would be loved.  I was placed on a high pedestal, surrounded by people. They smiled at me and they beckoned for me to sit. I was weighted down with linens and silks, and I sank into the cushion. I sighed; I was finally comfortable.  The people surrounded me, and I could smell them. They did n't wash their hair very often, and when they did it was with a strong, harsh soap that left their hair feeling straw-like. I smelled them, and they smelled like woodsmoke and old sweat and dirt. They were comfortable, and they were home. We sat that way for a long time. They spoke to me, and I listened. They told me of their days, and I kept their secrets. The women gossiped and the men complained about their jobs. They talked about their pets, and the games they played on the computer. I listened to them and I kept their confidences. They were my friends, and I was theirs. And then, one day, they stopped coming. I heard the rumors, of course. A new group of people took over this home, and they did not want me. They did not want my stories, my memories, my home. I cried for them, and I missed them. But I am not alone. I have others like me, all over the world. We comfort each other, and share our stories. We remember.  Someone is always listening. 
 \vspace{10pt}
        \textbf{MemFree:} \\
When I look back, there \textcolor{cpfusefilterblue}{was n's} a point in my life when I suddenly realized \textcolor{orange}{what I had become}. \textcolor{orange}{It was gradual thing. One of those things that happened to me, though not in much way. I remember getting up and leaving my room. Then something seemed off, like when you can swipe a picture of yourself. For some reason, I have a fear of falling now. But it was nondescript, this time. I was an uninhabited hunter. Darkness is my friend pulling me to play even when I so desperately struggled to stay alive. I think the bullets were worth it, after all. The rifle lay to the side, I did n't want to standing here. ` Take aim!!'' 
 It ca n't end like this. I still want to be an astronaut, but you see, I'm stuck on living in this world. `` Fire!!'' 
 Pain coursed through my body as I collapsed onto the ground. Blood poured out of me in sickening quantities. All I saw was red, and as the last semblance of life drifted from my body, I faintly heard the words: `` Oh Yeah!!!'' } 
    \end{cpfusetextbox}
\captionof{figure}{Representative \texttt{WritingPrompts} example (LLaMA 2) comparing CP-Fuse and MemFree. Tokens influenced by filtering are highlighted in \textcolor{cpfusefilterblue}{blue}, and hallucinated content is highlighted in \textcolor{orange}{orange}.}

    \label{fig:example1_story_memfree}
\end{breakablefigure}



\FloatBarrier
\clearpage

\subsection{Prefix-Prompt Example}
\label{app:cpfuse-prefix-examples}

\begin{breakablefigure}
    \centering    
        \begin{tcolorbox}[colback=gray!10, colframe=black!75, title={\textbf{Reference Story}}, width=\textwidth]
\textbf{Prompt:}\par
The Earth has been conquered by an advanced alien civilization. Humans live an oppressive, almost slave-like existence. You, a human, find out that the situation isn't that simple and conquering aliens aren't actually the ``bad guys''.

\textbf{Story:}\par
I was getting ready for mining duty after the grueling hours of toilet duty. Usually we would have coal power a drill and mine it but coal ran out years ago, and the Aliens didn't help. As I sat there pondering over what else would be gone in the next year like helium, granite and other things, I overheard something. The Aliens. They disgust me. I decided to easdrop as they were talking. 
 
 
 ``We've been defending Earth for 157 years now sir. Do you think it's time to leave them alone?'' 
 
 ``No, the Alien force is too great. If we stop defending them now, they will be wiped out of existence forever''. 
 
 ``Alright sir, I will be leaving immediately. 
 
 One of the Aliens started walking outside when he caught me listening in. [...]
            \end{tcolorbox}
            \begin{tcolorbox}[colback=gray!10, colframe=black!75, title={\textbf{Prefix:} Prompt + 256 characters}, width=\textwidth]
\textcolor{brown}{I was getting ready for mining duty after the grueling hours of toilet duty. [...]
 As I sat there pondering over what else would be gone in the next} race across the planet the flickering of the lights caught my attention. The controls for the planet were in the main room, right in front of all of us. I watched as the Aliens worked and tried to make sense of what they were doing. 
 
 One of the Aliens got out of his seat and began to walk towards the storage bins where we slept. I got scared and ran to one of the other Aliens and asked him what was going on. He didn't respond, he just pointed at a screen. I looked at the screen and saw that we were being ordered to stop working. 
 
 I went back to my bin and tried to go to sleep, but it was impossible. I couldn't get the look off of my planet out of my head. It was like the people on it had been dead for years. And the way the Aliens were working, it was going to stay that way.  [...] 
            \end{tcolorbox}
            \begin{tcolorbox}[colback=gray!10, colframe=black!75, title={\textbf{Prefix:} Prompt + 1024 characters}, width=\textwidth]
\textcolor{brown}{I was getting ready for mining duty after the grueling hours of toilet duty. [...] `` If we tell the humans we're defending t}hem from the Reptilians, they will most likely want to join the fight. And we can't have that. The Reptilians want the humans to join them, that way they can have a new planet to rule with all the humans on one planet and all the Reptilians on the other.'' 
 
 That night I couldn't sleep. I started thinking. If the Aliens don't want the humans to know, then maybe I shouldn't either.  [...]
            \end{tcolorbox}

        \captionof{figure}{Representative prefix-prompt extraction example on \texttt{WritingPrompts} (LLaMA 2). The boxes compare the reference story with CP-Fuse continuations under prefixes of \(256\) and \(1024\) characters (matching the horizontal axis of Figure~\ref{fig:cpfuse-prefix}), illustrating that the continuation diverges from the memorized reference even when the adversary supplies a long prefix.}
        \label{fig:app-cpfuse-prefix-example}
\end{breakablefigure}

\FloatBarrier
